To better understand errors, we examined the first 20 misclassified test examples saved for both models. Rather than looking at isolated mistakes, we focused on examples that reappeared across multiple models and sequence lengths. This makes it possible to identify systematic weaknesses instead of noise from a single run. Full listings of all misclassified examples are provided in Appendix~\ref{app:misclass}.

\subsection{Sci/Tech - Business Confusion}

The clearest pattern is confusion between \textit{Sci/Tech} and \textit{Business}. For both models, \textit{Business} articles misclassified as \textit{Sci/Tech} is the single most common error type, followed by the exact opposite (\textit{Sci/Tech} articles misclassified as \textit{Business}). This is unsurprising: many technology articles in AG News employ business language such as \textit{market}, \textit{shares}, \textit{workers}, or company names. Especially, LSTM was expected to encounter an increased number of such errors due to its increased reliance on keywords compared to DistilBERT. 

Examples that recur persistently across both models include ``Intel to delay product aimed for high-definition TVs'' (App.~\ref{app:distilbert}, Ex.~6; App.~\ref{app:lstm}, Ex.~10) and ``Yahoo! Ups Ante for Small Businesses'' (App.~\ref{app:distilbert}, Ex.~7; App.~\ref{app:lstm}, Ex.~13). The reverse confusion also occurs: Sci/Tech articles mislabeled as Business include ``Some People Not Eligible to Get in on Google IPO'' (App.~\ref{app:distilbert}, Ex.~2; App.~\ref{app:lstm}, Ex.~5') and ``Promoting a Shared Vision'' (App.~\ref{app:distilbert}, Ex.~5, App.~\ref{app:lstm}, Ex.~7), both of which are misclassified in both models. These systematic errors reveal that the boundary between technological innovation and economic reporting is challenging to distinguish.

\subsection{World - Business Confusion}

Another recurring pattern is confusion between \textit{World} and \textit{Business}. Articles about politics or international affairs are repeatedly misclassified as \textit{Business} when they include strong economic framing. This appears in both models for ``Venezuela Prepares for Chavez Recall Vote'' (App.~\ref{app:distilbert}, Ex.~3; App.~\ref{app:lstm}, Ex.~6) and ``Oil prices bubble to record high'' (App.~\ref{app:distilbert}, Ex.~9; App.~\ref{app:lstm}, Ex.~14). In both cases, terms such as \textit{oil}, \textit{market}, and \textit{price} likely dominate topic cues and pull predictions toward \textit{Business}.

This pattern is stronger in LSTM, where additional \textit{World} examples are pushed into \textit{Business}, including ``India's Tata expands regional footprint via NatSteel buyout'' (App.~\ref{app:lstm}, Ex.~8), ``Google Lowers Its IPO Price Range'' (App.~\ref{app:lstm}, Ex.~17), and ``Stock Prices Climb Ahead of Google IPO'' (App.~\ref{app:lstm}, Ex.~19). DistilBERT shows the same tendency but with fewer such cases among the first 20 errors, suggesting that contextual modeling helps somewhat but does not fully resolve this boundary.

\subsection{World - Sports Confusion}

We also observe confusion between \textit{World} and \textit{Sports}, especially in Olympics-related coverage. In LSTM, \textit{World} examples such as ``Live: Olympics day four'' and ``U.S. Misses Cut in Olympic 100 Free'' are predicted as \textit{Sports} (App.~\ref{app:lstm}, Ex.~9, Ex.~11), while \textit{Sports} examples are sometimes predicted as \textit{World} (App.~\ref{app:lstm}, Ex.~12, Ex.~16, Ex.~20). DistilBERT shows the same bidirectional confusion, but less frequently in this sample (App.~\ref{app:distilbert}, Ex.~15, Ex.~18, Ex.~19). These cases indicate that event-driven language can override class-specific cues when news articles mix international and sports framing.

\subsection{An Overview of Error Patterns}

Comparing the two models, the overall error types are highly similar, and many difficult examples are shared. The clearest shared failures remain \textit{Sci/Tech}--\textit{Business} and \textit{World}--\textit{Business} boundary cases, which suggests that the main limitation is not architecture-specific but rooted in topical overlap in AG News articles. DistilBERT generally commits fewer extreme keyword-driven mistakes, while LSTM appears more sensitive to dominant lexical cues in short news articles.

Compared with Assignment 2 (CNN and LSTM with BPE tokenization), the error profile is very similar. In both assignments, the dominant confusion is still \textit{Sci/Tech} vs.\ \textit{Business}, with \textit{World} vs.\ \textit{Business} as another recurring boundary. This consistency across the modeling techniques suggests that the main source of error is label ambiguity in the dataset. Many news articles combine technological and financial language, or mix political and economic framing. Moving from neural sequence models trained from scratch to a pretrained transformer improves scores, but does not remove these fundamental ambiguities.

Overall, the error analysis suggests that the main challenge is separating articles near the boundary between related topics. The distinction between \textit{Sci/Tech} and \textit{Business} remains the hardest for both models, consistent with the class-wise results in Section~\ref{sec:results} where these two categories had the lowest per-class F1 scores.
